\scanpage{75}
\mspage{81}

L'équiv. $a\Longleftrightarrow b')$ \deleted{\unsure{se}} se déduit de $a)\Longrightarrow b)$ par dualité (la condition a) étant \unsure{autoduale}), ou en reprenant le $\deleted{\unsure{\Leftrightarrow}}=$ raisonnement pour $a\Rightarrow b'$, et pour $b')\Rightarrow a)$ aussi (en \deleted{changeant} remplaçant la coh. par l'homologie).

\uline{Définition 2} \unsure{Si} \deleted{\unsure{de}} les conditions équivalentes de la prop. 2 sont satisfaites, on dit que $f$ est une \uline{$\mathbb D$-fibration}.

\margin{\uline{NB} C'est une notion à \uline{la fois} locale et ``colocale'' sur $Y$ (alors que fibr. g. est locale, fibr. dr. est colocale)}

\uline{Proposition 3} a) La notion de $\mathbb D$-fibration est stable par changt de base \uline{et par composition}\textsuperscript{*}

% slanted
\margin{?? Stabilité par composition \uline{n'est} pas claire}

b) Soient
\diagram{\begin{tikzcd}[ampersand replacement=\&]
X\arrow[rr,"f"]\arrow[dr,dash,"p"'] \& {} \& X'\arrow[dl,dash,"p'"]\\
\& Y\arrow[u,phantom,"\text{comm}" description] \&
\end{tikzcd}}
avec $p,p'$ des $\mathbb D$-fibrations \inserted{(il \unsure{suffira|suffirait} que $p$ et $p'$ soient des $\mathbb D$-fibr. coh., ou que ce soient des $\mathbb D$-fibr. homologiques, ou que $p,p'$ soient propres, ou (t.t.\unsure{lisses|lisse}))$;$} \deleted{\unsure{pour que}} $f$ induit des $\mathbb D$-équiv. sur les fibres \unsure{soit}; alors $f$ est une $\mathbb D$-équivalence.

% slanted
\margin{Peut se préciser considérablement dans le cas \deleted{\unsure{où}} où $p,p'$ sont des $\mathbb D$-fibrations. Cf prop. 6 page 84}

\uline{Dém} a) est immédiat sur la condition a) (prop 2 :) de la notion de $\mathbb D$-fibration. Pour b), supposons \inserted{que} $f$ induise des $\mathbb D$-équiv. sur les fibres. On veut prouver que
\[
H^{\unsure{\bullet|\cdot|\prime}}_{\mathbb D}(X,M)\unsure{-|\longrightarrow} H^\bullet_{\mathbb D}(X',M)
\]
est un iso, pour tt $M$ dans $A$, mais cette flèche s'identifie à
\[
H^\bullet_{\mathbb D}(Y,f_*(M_X))\longrightarrow H^\bullet_{\mathbb D}(Y,f'_*(M_{X'}))
\]
déduite de
\[
f_*(M_X)\unsure{-|\longrightarrow} f_*(M_{X'})\qquad(\mathrm{NB}\quad M_X\simeq f^*M_{X'}),
\]
% archive sheet number: 74
